\documentclass{article}
\usepackage[T1]{fontenc}
\usepackage{lmodern}
\usepackage{graphicx}
\usepackage{amsmath,amssymb}
\usepackage[a4paper,margin=2cm]{geometry}
\usepackage[absolute,overlay]{textpos}
\setlength{\TPHorizModule}{1mm}
\setlength{\TPVertModule}{1mm}
\usepackage[indonesian]{babel}
\usepackage{hyperref}
\setlength{\emergencystretch}{2em}
% unit: Halaman 20

\begin{document}

% === Halaman 20 (llm) ===

\begin{itemize}
    \item Algoritma Diffie-Hellman dipakai di dalam protokol SSL dan TLS. SSL/TLS digunakan untuk mengamankan komunikasi di internet antara \textit{client} dan \textit{server}.
    \item Pada tahun 2002, Hellman menyarankan algoritma DH diberi nama algoritma pertukaran kunci \textbf{Diffie--Hellman--Merkle} sebagai penghargaan terhadap kontribusi Ralph Merkle dalam penemuan kriptografi kunci-publik:
    \item Hellman menulis:
    \begin{quote}
    \textit{The system...has since become known as Diffie--Hellman key exchange. While that system was first described in a paper by Diffie and me, it is a public key distribution system, a concept developed by Merkle, and hence should be called 'Diffie--Hellman--Merkle key exchange' if names are to be associated with it. I hope this small pulpit might help in that endeavor to recognize Merkle's equal contribution to the invention of public key cryptography.[7]}
    \end{quote}
\end{itemize}

\begin{flushright}
Sumber: Wikipedia
\end{flushright}

\end{document}
